\documentclass[11pt]{article}
\usepackage{amsmath,amssymb,amsthm}
\usepackage[margin=1in]{geometry}

\newtheorem{theorem}{Theorem}
\newtheorem{lemma}{Lemma}
\newtheorem{corollary}{Corollary}
\newcommand{\bits}{\{0,1\}}

\title{Laboratory V95: Balanced Canonical Loaders and a PP-Hard Next Bit}
\author{NC0\_k-Avoid Laboratory}
\date{2026-08-10}

\begin{document}
\maketitle

Let $C:\bits^N\to\bits^M$.  The canonical rule maintains a fiber $S_j$ and
chooses output zero at position $j$ exactly when the zero-child is no larger
than the one-child, breaking ties toward zero.

\begin{lemma}[Balanced definition]
Let $z$ be an input coordinate absent from all previous output functions and
let
\[
 D_f(z,a,b)=z\oplus f(a,b)
\]
for any binary Boolean function $f$.  Under every prior prefix, the zero and
one children of $D_f$ have exactly equal size.  Hence the canonical bit is zero,
and conditioning on it uniquely imposes $z=f(a,b)$.
\end{lemma}

\begin{proof}
Toggling only $z$ preserves all previous outputs and flips $D_f$.  This is a
fixed-point-free involution between the two children.  The tie-zero rule chooses
zero, where $z=f(a,b)$ holds uniquely.
\end{proof}

\begin{corollary}[Composable loader]
Any acyclic fan-in-two Boolean computation can be installed by introducing one
fresh input coordinate per internal wire and listing the corresponding balanced
definition outputs in topological order.  The entire loader prefix is
canonically all zero, and every source assignment has one extension encoding
all correct wire values.
\end{corollary}

Now take nonempty 3-CNFs $F_0,F_1$ over the same $n$ variables.  Compute each
clause using two balanced signed-OR definitions and combine clauses by balanced
AND definitions.  Let the resulting truth coordinates be $t_0,t_1$, let $s$ be
one free selector coordinate, and let $q$ be the number of definition outputs.
After the genuine canonical prefix $0^q$, there is one extension for each
$(x,s)$.

Place next the ternary output
\[
 H(s,t_0,t_1)=
 \begin{cases}
  \neg t_0,&s=0,\\
  t_1,&s=1.
 \end{cases}
\]
If $A=\#\mathrm{sat}(F_0)$ and $B=\#\mathrm{sat}(F_1)$, then
\[
 N_C(0^q0)=A+(2^n-B),\qquad
 N_C(0^q1)=(2^n-A)+B.
\]
Consequently
\[
 N_C(0^q0)\le N_C(0^q1)
 \quad\Longleftrightarrow\quad
 A\le B.
\]

\begin{theorem}[PP-hard genuine canonical next bit]
There is a fixed seven-type output language of locality at most three such that
computing a specified next bit of the canonical halving sequence is PP-hard for
$\mathrm{NC}^0_3$ circuits of exact stretch $M=N+1$, even when every preceding
canonical decision is an exact tie and the preceding canonical prefix is
syntactically $0^q$.
\end{theorem}

\begin{proof}
V94 establishes PP-completeness of comparing the numbers of satisfying
assignments of two 3-CNFs.  The balanced loader above makes the prefix genuinely
canonical, and the displayed child-count identity makes the next canonical bit
zero exactly for the non-strict source comparison.  The compiler uses four
signed-OR balanced-definition tables, one AND balanced-definition table, $H$,
and projection padding.  If the formulas use $r_0,r_1$ clauses, then
$q=3(r_0+r_1)-2$ and $N=n+1+q$.  Appending exactly $n+1$ projection outputs
gives $M=N+1$.
\end{proof}

\begin{corollary}
Exact computation of the complete canonical halving word for all such circuits
is PP-hard.  A deterministic polynomial-time algorithm for that specific word
would imply $P=PP$ and hence $P=NP$.
\end{corollary}

\paragraph{Nonclaim.}
This does not prove that finding an arbitrary avoided output is PP-hard.  Range
avoidance is a total search problem with many possible correct outputs; V95 only
closes the strategy of reproducing this exact canonical word.

\end{document}
